\chapter{Anomaly Inflow and Invertible Field Theories}
\label{ch:global-anomaly-inflow-invertible-field-theories}

\ConventionsMixed

An anomaly of a \(D\)-dimensional QFT with background fields is an obstruction
to making the background-field dependence an ordinary function on the
background-field moduli stack.  Inflow records that obstruction as the
boundary variation of a \((D+1)\)-dimensional invertible field theory.

\section{Background Fields and Relative Partition Functions}

Let \(M^D\) be a closed spacetime and let \(A\) denote a background field for a
global symmetry, including possible higher-form components.  Write
\[
  \mathcal B(M)
\]
for the groupoid of such backgrounds on \(M\).  A nonanomalous partition
function assigns a complex number \(Z_M(A)\) invariant under isomorphisms in
\(\mathcal B(M)\).  An anomalous theory assigns instead an element
\[
  Z_M(A)\in \mathcal L_M(A)
\]
of a one-dimensional complex vector space.  The collection \(\mathcal L_M\)
over \(\mathcal B(M)\) is the anomaly line.

\section{Invertible Bulk Theory}

A \((D+1)\)-dimensional invertible theory assigns a line
\[
  \mathcal Z_N(\widetilde A)
\]
to a closed \((D+1)\)-manifold \(N\) with background \(\widetilde A\), and a
linear map to bordisms.  When the bulk is represented by a local action
\[
  \exp\!\left(2\pi\ii\int_N I_{D+1}(\widetilde A)\right),
\]
the integrand \(I_{D+1}\) is a cocycle or differential cocycle whose periods
are integral in the normalization used.  If \(N\) has boundary \(M\), the same
bulk theory assigns a boundary line to \(A=\widetilde A|_M\).  A relative
\(D\)-dimensional QFT is a boundary condition for this invertible theory, and
its partition function is a vector in the boundary line.

\section{Descent for Connected Symmetries}

For a continuous symmetry with connection \(A\), let \(P_{D+2}(F)\) be an
invariant polynomial in the curvature \(F=dA+A^2\) with integral periods.
Choose a Chern-Simons form \(I_{D+1}(A)\) satisfying
\[
  dI_{D+1}(A)=P_{D+2}(F).
\]
For an infinitesimal gauge parameter \(\alpha\),
\[
  \delta_\alpha I_{D+1}(A)=dI_D^{(1)}(\alpha,A).
\]
Therefore on a manifold \(N\) with boundary \(M\),
\[
  \delta_\alpha
  2\pi\ii\int_N I_{D+1}
  =
  2\pi\ii\int_M I_D^{(1)}(\alpha,A).
\]
The boundary generating functional must vary by the negative of this term for
the combined bulk-boundary partition function to be gauge invariant.

\begin{proposition}[Cocycle ambiguity and local counterterms]
\label{prop:inflow-cocycle-counterterm}
Changing \(I_{D+1}\) by \(I_{D+1}+dK_D\) changes the boundary anomaly by the
variation of the local \(D\)-dimensional functional
\[
  2\pi\ii\int_M K_D(A).
\]
Thus the anomaly invariant is the cohomology class of the bulk invertible
theory, while a representative anomaly current depends on a boundary
counterterm coordinate.
\end{proposition}

\begin{proof}
On a manifold with boundary,
\[
  \int_N dK_D=\int_M K_D .
\]
Under a background gauge transformation, the additional boundary variation is
exactly the variation of this \(D\)-dimensional local functional.  Hence two
bulk cocycles differing by a coboundary define anomaly representatives that
are related by a local counterterm.
\end{proof}

\section{One-Form Symmetry Example}

Let a four-dimensional theory have a finite abelian one-form symmetry \(A\).
A background is a two-cocycle
\[
  B\in Z^2(M,A).
\]
An anomaly involving this symmetry is represented by a five-dimensional
cohomology class
\[
  \omega\in H^5(B^2A,U(1)).
\]
For \(A=\mathbb Z_N\), expressions built from cup products and the Pontryagin
square give concrete representatives after the manifold orientation and spin
structure have been specified.  The boundary theory then has line operators
whose fusion and braiding with surface defects are acted on by the anomaly
line; changing a four-dimensional discrete theta term shifts the
representative but not the inflow class unless the bulk cohomology class
changes.

\section{Gauging Criterion}

To gauge a global symmetry, one must sum or integrate over its backgrounds.
For an anomalous theory this operation is possible only after the anomaly
line is trivialized, or after the theory is coupled to another system whose
anomaly line is the inverse.  In the invertible-language formulation the
criterion is
\[
  \mathcal Z_{\mathrm{bulk}}\otimes
  \mathcal Z_{\mathrm{counter}}
  \simeq 1
\]
as an invertible theory over the background-field bordism category.  A local
counterterm supplies such a trivialization exactly when the anomaly class is
a coboundary in the chosen background-field cohomology theory.

\begin{openproblem}[Operator-level inflow]
\label{op:operator-level-inflow}
Develop a reconstruction theorem that starts from an anomalous local QFT,
constructs its anomaly line over background fields, and derives the action of
the inflow theory on extended operators, defect junctions, and Hilbert spaces
on manifolds with boundary.
\end{openproblem}
